\documentclass[12pt,aspectratio=169]{beamer}
\usepackage[utf8]{inputenc}
\usepackage[T2A]{fontenc}
\usepackage[main=russian]{babel}
\usepackage{amsmath}
\usepackage{amsfonts}
\usepackage{array}
\usepackage{booktabs}

\usetheme{default}
\usecolortheme{whale}
\setbeamertemplate{navigation symbols}{}
\setbeamertemplate{footline}[page number]

\setlength{\parindent}{0pt}

\title{Метод Жордана-Гаусса}
\subtitle{Решение систем линейных алгебраических уравнений}
\author{В.\,Ю.\,Винник}
\institute{Гимназия №6 г.\,Минска}
\date{8-Г класс}

\begin{document}

\begin{frame}
    \titlepage
\end{frame}

\begin{frame}{Решим систему двух уравнений}
\textbf{Дано:}
        \[
            \left\{
            \begin{array}{ccccc}
              x  & - & y & = & 2 \\
              2x & + & y & = & 7
            \end{array}
            \right.
        \]
        \onslide<2->{
Выразим $x$ через $y$ из первого уравнения:
            \[
                x = y + 2
            \]
        }
        \onslide<3->{
Подставим правую часть вместо $x$ во второе уравнение:
            \begin{align*}
                2(y+2) + y &= 7 \\
                2y + 4 + y &= 7 \\
                3y &= 3 \\
                y &= 1
            \end{align*}
        }
        \onslide<4->{
Подставляем найденное значение $y$ и находим $x$:
            \[
                x = 1 + 2 = 3
            \]
        }
\end{frame}

\begin{frame}{А если уравнений три?}
        \[
            \left\{
            \begin{array}{ccccccc}
              x  & + &  y & + &  z & = &  8 \\
              x  & + & 2y & + & 4z & = & 16 \\
             3x  & + &  y & + & 2z & = & 13
            \end{array}
            \right.
        \]
        \onslide<2->{
            Выразим $z$ через $x$ и $y$ из первого уравнения:
            \[
                z = 8 - x - y
            \]
        }
        \onslide<3->{
            Подставим во второе уравнение:
            \[
                \begin{array}{rcl}
                    x + 2y + 4(8 - x - y) & = & 16 \\
                    x + 2y + 32 - 4x - 4y & = & 16 \\
                  -3x - 2y & = & -16 \\
                   3x + 2y & = & 16
                \end{array}
            \]
        }
\end{frame}

\begin{frame}{Три уравнения --- продолжение}
    Подставим $z = 8 - x - y$ в третье уравнение:
    \[
        \begin{array}{rcl}
             3x + y + 2z = 13 \\
             3x + y + 2(8 - x - y) = 13 \\
             3x + y + 16 - 2x - 2y = 13 \\
             x - y = -3
        \end{array}
    \]
    \onslide<2->{
        Получили систему из двух уравнений с двумя неизвестными:
        \[
            \left\{
            \begin{array}{ccccc}
              3x & + & 2y & = & 16 \\
               x & - &  y & = & -3
            \end{array}
            \right.
        \]
        Её можно решить подстановкой, выразив $y$ через $x$ (или наоборот), как в первом примере.
    }
\end{frame}

\begin{frame}{Обобщение примеров}
    \begin{itemize}
        \item<1-> Системы из трёх уравнений с тремя неизвестными подстановкой решать можно.
        \item<2-> Таким же способом можно решать системы и с большим числом неизвестных и уравнений.
        \item<3-> Однако решение получается громоздким и неудобным.
        \item<4-> Нужно делать много промежуточных выкладок.
        \item<5-> В них легко запутаться, а найти ошибку трудно.
        \item<6-> В реальных задачах (проектирование технических устройств, управление экономикой) бывают десятки, сотни, тысячи неизвестных и уравнений!
        \item<7-> Нужен удобный, простой и универсальный метод!
    \end{itemize}
\end{frame}

\begin{frame}[fragile]{Терминология}
    \begin{enumerate}
        \item \textbf{СЛАУ} --- \textbf{с}истема \textbf{л}инейных \textbf{а}лгебраических \textbf{у}равнений.
        \item \textbf{Алгебраические} уравнения (в отличие от тригонометрических, дифференциальных и др.) --- содержат только операции сложения, вычитания, умножения и деления.
        \item \textbf{Линейные} --- неизвестные входят только в первой степени (не в квадрате, кубе или более высоких степенях). Неизвестные не перемножаются между собой, только умножаются на постоянные коэффициенты.
    \end{enumerate}
    В общем случае СЛАУ имеет вид
    \[
    \begin{array}{cccccccc}
        a_{11}x_1 & + & a_{12}x_2 & + & \ldots & + a_{1n}x_n & = & b_1 \\
        a_{21}x_1 & + & a_{22}x_2 & + & \ldots & + a_{2n}x_n & = & b_2 \\
        \ldots \\
        a_{n1}x_1 & + & a_{n2}x_2 & + & \ldots & + a_{nn}x_n & = & b_n
    \end{array}
    \]
\end{frame}

\begin{frame}{Убираем лишнее}
    \begin{block}{Достаточно таблицы коэффициентов}
        Вместо записи со всеми \(x\), \(+\), \(=\) можно записывать только числа:
        \[
        \left(\begin{array}{cccc|c}
            a_{11} & a_{12} & \ldots & a_{1n} & b_1 \\
            a_{21} & a_{22} & \ldots & a_{2n} & b_2 \\
            \vdots & \vdots & \ddots & \vdots & \vdots \\
            a_{n1} & a_{n2} & \ldots & a_{nn} & b_n
        \end{array}\right)
        \]
    \end{block}
\end{frame}

\begin{frame}{Пример}
   Система уравнений
    \[
        \left\{
        \begin{array}{ccccccc}
          x_1 & + &  x_2 & + &  x_3 & = &  8 \\
          x_1 & + & 2x_2 & + & 4x_3 & = & 16 \\
         3x_1 & + &  x_2 & + & 2x_3 & = & 13
        \end{array}
        \right.
    \]
    \onslide<2->{
            Превращается в таблицу
            \[
            \left(\begin{array}{ccc|c}
                1 & 1 & 1 & 8 \\
                1 & 2 & 4 & 16 \\
                3 & 1 & 2 & 13
            \end{array}\right)
            \]
    }
\end{frame}

\begin{frame}{Свойства СЛАУ}
    \begin{block}{Свойство 1}
        Если обе части уравнения умножить на одно и то же число \(q \neq 0\), то решение \textbf{не изменится}.
    \end{block}
    \onslide<2->{
        \[
        a_{i1}x_1 + a_{i2}x_2 + \ldots + a_{in}x_n = b_i
        \]
        \vspace{-0.3cm}
        \[
        \Big\Updownarrow
        \]
        \[
        q \cdot a_{i1}x_1 + q \cdot a_{i2}x_2 + \ldots + q \cdot a_{in}x_n = q \cdot b_i
        \]
    }
    \onslide<3->{
        Любую строку в таблице можно умножить на число (не равное нулю).
    }
\end{frame}

\begin{frame}{Свойства СЛАУ --- продолжение}
    \begin{block}{Свойство 2}
        Если к одному уравнению прибавить другое (почленно), то решение \textbf{не изменится}.
    \end{block}
    \onslide<2->{
        \[
        \begin{array}{rcl}
            a_{i1}x_1 + \ldots + a_{in}x_n &=& b_i \\
            a_{j1}x_1 + \ldots + a_{jn}x_n &=& b_j
        \end{array}
        \]
        \vspace{-0.3cm}
        \[
        \Big\Downarrow
        \]
        \[
        (a_{i1} + a_{j1})x_1 + \ldots + (a_{in} + a_{jn})x_n = b_i + b_j
        \]
    }
    \onslide<3->{
        К любой строке таблицы можно прибавить (поэлементно) любую другую.
    }
\end{frame}

\begin{frame}{Совместное применение свойств 1 и 2}
    Пусть даны два уравнения:
    \[
    \left\{
    \begin{array}{rrrrr}
        2x &+& 3y &=& 8 \\
         x &-&  y &=& 1
    \end{array}
    \right.
    \]
    \onslide<2->{
        \vspace{0.3cm}
        Умножим второе на \(q=2\) и прибавим к первому:
        \[
        \begin{array}{rrrrrr}
              & 2x &+& 3y &=& 8 \\
            + & 2x &-& 2y &=& 2 \\
            \hline
              & 4x &+& y &=& 10
        \end{array}
        \]
    }
    \onslide<3->{
            Система
            \[
            \left\{
            \begin{array}{rrrrr}
                4x &+& y &=& 10 \\
                 x &-& y &=& 1
            \end{array}
            \right.
            \]
            имеет \textbf{то же решение}, что и исходная!
    }
\end{frame}

\begin{frame}{То же самое в виде таблиц}
    Исходная таблица
    \[
    \left(\begin{array}{cc|c}
        2 & 3 & 8 \\
        1 & -1 & 1
    \end{array}\right)
    \]
    \onslide<2->{
        Вторую строку умножить на 2 и прибавить к первой:
        \[
        \begin{array}{ccc|c}
               & 2 &  3 & 8 \\
            +  & 2 & -2 & 2 \\
            \hline
               & 4 &  1 & 10
        \end{array}
        \]
    }
    \onslide<2->{
        Преобразованная таблица
        \[
        \left(\begin{array}{cc|c}
            4 & 1 & 10 \\
            1 & -1 & 1
        \end{array}\right)
        \]
    }
\end{frame}

\begin{frame}{Главная идея}
    \begin{center}
        \Large
        С помощью описанных свойств менять вид системы, \\
        чтобы решение СЛАУ осталось неизменным, \\
        но стало очевидным!
    \end{center}
    \onslide<2->{
        \begin{itemize}
            \item Исключим \(x_1\) из всех уравнений, кроме первого.
            \item То есть превратим коэффициенты при $x_1$ в нули.
            \item Затем \(x_2\) исключим из всех уравнений, кроме первого и второго.
            \item И так далее...
        \end{itemize}
    }
\end{frame}

\begin{frame}{Пример: система 3x3 (шаг 1)}
    Вернёмся к системе, которую решали подстановкой:
    \[
    \left\{
    \begin{array}{ccccccc}
      x  & + &  y & + &  z & = &  8 \\
      x  & + & 2y & + & 4z & = & 16 \\
     3x  & + &  y & + & 2z & = & 13
    \end{array}
    \right.
    \]
    \vspace{0.3cm}
    \onslide<2->{
        Запишем её в виде таблицы:
        \[
        \left(\begin{array}{ccc|c}
            1 & 1 & 1 & 8 \\
            1 & 2 & 4 & 16 \\
            3 & 1 & 2 & 13
        \end{array}\right)
        \]
    }
    \onslide<3->{
        \textbf{Цель прямого хода:} получить нули под главной диагональю.
    }
\end{frame}

\begin{frame}{Пример 3x3: исключаем \(x\) из 2-й строки}
    \[
    \left(\begin{array}{ccc|c}
        1 & 1 & 1 & 8 \\
        1 & 2 & 4 & 16 \\
        3 & 1 & 2 & 13
    \end{array}\right)
    \]
    \vspace{0.3cm}
    \onslide<2->{
        Умножим 1-ю строку на \(-1\) и прибавим ко 2-й:
        \[
        \begin{array}{ccc|c}
            1 & 2 & 4 & 16 \\
            -1 & -1 & -1 & -8 \\
            \hline
            0 & 1 & 3 & 8
        \end{array}
        \]
    }
    \onslide<3->{
        Получаем таблицу с нулём в позиции \((2,1)\):
        \[
        \left(\begin{array}{ccc|c}
            1 & 1 & 1 & 8 \\
            0 & 1 & 3 & 8 \\
            3 & 1 & 2 & 13
        \end{array}\right)
        \]
    }
\end{frame}

\begin{frame}{Пример 3x3: исключаем \(x\) из 3-й строки}
    \[
    \left(\begin{array}{ccc|c}
        1 & 1 & 1 & 8 \\
        0 & 1 & 3 & 8 \\
        3 & 1 & 2 & 13
    \end{array}\right)
    \]
    \vspace{0.3cm}
    \onslide<2->{
        Умножим 1-ю строку на \(-3\) и прибавим к 3-й:
        \[
        \begin{array}{ccc|c}
            3 & 1 & 2 & 13 \\
            -3 & -3 & -3 & -24 \\
            \hline
            0 & -2 & -1 & -11
        \end{array}
        \]
    }
    \onslide<3->{
        Получаем таблицу с нулями в первом столбце (кроме строки 1):
        \[
        \left(\begin{array}{ccc|c}
            1 & 1 & 1 & 8 \\
            0 & 1 & 3 & 8 \\
            0 & -2 & -1 & -11
        \end{array}\right)
        \]
    }
\end{frame}

\begin{frame}{Пример 3x3: исключаем \(y\) из 3-й строки}
    \[
    \left(\begin{array}{ccc|c}
        1 & 1 & 1 & 8 \\
        0 & 1 & 3 & 8 \\
        0 & -2 & -1 & -11
    \end{array}\right)
    \]
    \vspace{0.3cm}
    \onslide<2->{
        Умножим 2-ю строку на \(2\) и прибавим к 3-й:
        \[
        \begin{array}{ccc|c}
            0 & -2 & -1 & -11 \\
            0 & 2 & 6 & 16 \\
            \hline
            0 & 0 & 5 & 5
        \end{array}
        \]
    }
    \onslide<3->{
        \textbf{Прямой ход завершён!} Таблица приведена к верхнетреугольному виду
        \[
        \left(\begin{array}{ccc|c}
            1 & 1 & 1 & 8 \\
            0 & 1 & 3 & 8 \\
            0 & 0 & 5 & 5
        \end{array}\right)
        \]
    }
\end{frame}

\begin{frame}{Пример 3x3: обратный ход, обнуляем над диагональю}
    После прямого хода имеем верхнетреугольную таблицу:
    \[
    \left(\begin{array}{ccc|c}
        1 & 1 & 1 & 8 \\
        0 & 1 & 3 & 8 \\
        0 & 0 & 5 & 5
    \end{array}\right)
    \]
    \onslide<2->{
        \textbf{Цель обратного хода:} получить слева единичную матрицу, то есть таблицу с одними лишь единицами на главной диагонали.
    }
\end{frame}

\begin{frame}{Пример 3x3: делаем диагональ единичной}
    \[
    \left(\begin{array}{ccc|c}
        1 & 1 & 1 & 8 \\
        0 & 1 & 3 & 8 \\
        0 & 0 & 5 & 5
    \end{array}\right)
    \]
    \vspace{0.3cm}
    \onslide<2->{
        Разделим 3-ю строку на \(5\):
        \[
        \left(\begin{array}{ccc|c}
            1 & 1 & 1 & 8 \\
            0 & 1 & 3 & 8 \\
            0 & 0 & 1 & 1
        \end{array}\right)
        \]
    }
    \onslide<3->{
        Теперь на главной диагонали стоят единицы.
    }
\end{frame}

\begin{frame}{Пример 3x3: обнуляем элемент \((2,3)\)}
    \[
    \left(\begin{array}{ccc|c}
        1 & 1 & 1 & 8 \\
        0 & 1 & 3 & 8 \\
        0 & 0 & 1 & 1
    \end{array}\right)
    \]
    \vspace{0.3cm}
    \onslide<2->{
        Обнулим элемент \((2,3)\). Умножим 3-ю строку на \(-3\) и прибавим ко 2-й:
        \[
        \begin{array}{ccc|c}
            0 & 1 & 3 & 8 \\
            0 & 0 & -3 & -3 \\
            \hline
            0 & 1 & 0 & 5
        \end{array}
        \]
    }
    \onslide<3->{
        \[
        \left(\begin{array}{ccc|c}
            1 & 1 & 1 & 8 \\
            0 & 1 & 0 & 5 \\
            0 & 0 & 1 & 1
        \end{array}\right)
        \]
    }
\end{frame}

\begin{frame}{Пример 3x3: обнуляем элемент \((1,3)\)}
    \[
    \left(\begin{array}{ccc|c}
        1 & 1 & 1 & 8 \\
        0 & 1 & 0 & 5 \\
        0 & 0 & 1 & 1
    \end{array}\right)
    \]
    \vspace{0.3cm}
    \onslide<2->{
        Обнулим элемент \((1,3)\). Умножим 3-ю строку на \(-1\) и прибавим к 1-й:
        \[
        \begin{array}{ccc|c}
            1 & 1 & 1 & 8 \\
            0 & 0 & -1 & -1 \\
            \hline
            1 & 1 & 0 & 7
        \end{array}
        \]
    }
    \onslide<3->{
        \[
        \left(\begin{array}{ccc|c}
            1 & 1 & 0 & 7 \\
            0 & 1 & 0 & 5 \\
            0 & 0 & 1 & 1
        \end{array}\right)
        \]
    }
\end{frame}

\begin{frame}{Пример 3x3: обнуляем элемент \((1,2)\)}
    \[
    \left(\begin{array}{ccc|c}
        1 & 1 & 0 & 7 \\
        0 & 1 & 0 & 5 \\
        0 & 0 & 1 & 1
    \end{array}\right)
    \]
    \vspace{0.3cm}
    \onslide<2->{
        Обнулим элемент \((1,2)\). Умножим 2-ю строку на \(-1\) и прибавим к 1-й:
        \[
        \begin{array}{ccc|c}
            1 & 1 & 0 & 7 \\
            0 & -1 & 0 & -5 \\
            \hline
            1 & 0 & 0 & 2
        \end{array}
        \]
    }
    \onslide<3->{
        \[
        \left(\begin{array}{ccc|c}
            1 & 0 & 0 & 2 \\
            0 & 1 & 0 & 5 \\
            0 & 0 & 1 & 1
        \end{array}\right)
        \]
    }
    \onslide<4->{
        \textbf{Получили единичную матрицу слева!}
    }
\end{frame}

\begin{frame}{Пример 3x3: ответ}
    \[
    \left(\begin{array}{ccc|c}
        1 & 0 & 0 & 2 \\
        0 & 1 & 0 & 5 \\
        0 & 0 & 1 & 1
    \end{array}\right)
    \]
    \vspace{0.3cm}
    \onslide<2->{
        Эта таблица соответствует системе:
        \[
        \left\{
        \begin{array}{ccccc}
            x &   &   & = & 2 \\
              & y &   & = & 5 \\
              &   & z & = & 1
        \end{array}
        \right.
        \]
    }
    \onslide<3->{
        \begin{block}{Ответ}
            \[
            \boxed{x=2,\quad y=5,\quad z=1}
            \]
        \end{block}
    }
\end{frame}

\begin{frame}{Пример 3x3: проверка}
    Подставим найденные значения в исходную систему:
    \[
    \left\{
    \begin{array}{ccccccc}
      x  & + &  y & + &  z & = &  8 \\
      x  & + & 2y & + & 4z & = & 16 \\
     3x  & + &  y & + & 2z & = & 13
    \end{array}
    \right.
    \]
    \vspace{0.3cm}
    \onslide<2->{
        \[
        2+5+1=8
        \]
        \[
        2+2\cdot 5+4\cdot 1=2+10+4=16
        \]
        \[
        3\cdot 2+5+2\cdot 1=6+5+2=13
        \]
    }
    \onslide<3->{
        \begin{exampleblock}{Вывод}
            Решение найдено верно.
        \end{exampleblock}
    }
\end{frame}

\begin{frame}{Алгоритм метода Жордана-Гаусса}
    \begin{enumerate}
        \item \textbf{Прямой ход:} приводим таблицу к \alert{верхнетреугольному} виду:
        \begin{itemize}
            \item последовательно исключаем \(x_1\), затем \(x_2\), затем \(x_3\) и так далее;
            \item в результате под главной диагональю получаются нули.
        \end{itemize}
        \item \textbf{Обратный ход:} преобразуем таблицу к \alert{единичной}:
        \begin{itemize}
            \item делим каждую строку на её диагональный элемент, чтобы на главной диагонали получились единицы;
            \item затем последовательно обнуляем элементы \alert{над} главной диагональю;
        \end{itemize}
        \item В результате таблица принимает вид
        \[
        \left(\begin{array}{cccc|c}
            1 & 0 & \ldots & 0 & x_1 \\
            0 & 1 & \ldots & 0 & x_2 \\
            \vdots & \vdots & \ddots & \vdots & \vdots \\
            0 & 0 & \ldots & 1 & x_n
        \end{array}\right),
        \]
        и ответ читается сразу: \(x_1,\ x_2,\ \ldots,\ x_n\).
    \end{enumerate}
\end{frame}

\begin{frame}{Особый случай: нули на главной диагонали}
    \[
    \left(\begin{array}{ccc|c}
        1 & 0 & 0 & 3 \\
        0 & 0 & 0 & 4 \\
        0 & 0 & 1 & 1
    \end{array}\right)
    \]
    \begin{itemize}
        \item<2-> На месте \(a_{22}\) стоит \alert{0}.
        \item<3-> Это соответствует уравнению $0\cdot y = 4$.
        \item<4-> Такое уравнение очевидно не имеет решений!
        \item<5-> Следовательно, решений не имеет и вся СЛАУ.
    \end{itemize}
\end{frame}

\begin{frame}{Особый случай --- продолжение}
    \[
    \left(\begin{array}{ccc|c}
        1 & 1 & 1 & 3 \\
        0 & 0 & 1 & 1 \\
        0 & 0 & 0 & 0
    \end{array}\right)
    \]
    \begin{itemize}
        \item<2-> В строке 3 нулевой не только коэффициент при неизвестной, но и свободный член.
        \item<3-> Это соответствует уравнению $0\cdot z = 0$.
        \item<4-> Это вырожденное уравнение, оно истинно при любом $z$.
        \item<5-> СЛАУ имеет бесконечно много решений.
        \item<6-> Метод Жордана-Гаусса не только находит решение, но и показывает, \textbf{существует ли} оно и \textbf{единственно} ли.
    \end{itemize}
\end{frame}

\begin{frame}{Итоги}
    \begin{itemize}
        \item \alert{СЛАУ} — системы линейных алгебраических уравнений.
        \item \alert{Метод Жордана-Гаусса} — универсальный способ решения.
        \item Два этапа:
        \begin{enumerate}
            \item \textbf{Прямой ход} $\rightarrow$ верхнетреугольный вид.
            \item \textbf{Обратный ход} $\rightarrow$ нахождение неизвестных.
        \end{enumerate}
        \item Разрешает:
        \begin{itemize}
            \item единственное решение;
            \item бесконечно много решений;
            \item отсутствие решений.
        \end{itemize}
        \item Удобен для реализации в виде программы для компьютера.
        \item Применяется в науке, технике, экономике для систем с \alert{тысячами} уравнений.
    \end{itemize}
    \vspace{0.5cm}
    \begin{center}
        \textbf{Спасибо за внимание!}
    \end{center}
\end{frame}

\end{document}
